\documentclass[11pt]{article}
\usepackage[margin=1in]{geometry}
\usepackage{amsmath,amssymb}
\usepackage{graphicx}
\usepackage{hyperref}
\usepackage{enumitem}
\usepackage{titlesec}
\usepackage{booktabs}

\titleformat{\section}{\large\bfseries}{\thesection}{1em}{}
\titleformat{\subsection}{\normalsize\bfseries}{\thesubsection}{1em}{}

\title{\textbf{Postdoctoral Research Proposal}\\[0.3em]
\Large Explainable AI for Human-AI Collaboration in Security Operations Centers}

\author{Research Project 4: Continuous-Time Explainability for SOC Analysts}
\date{}

\begin{document}

\maketitle

\section{Executive Summary}

This proposal outlines a 22-month postdoctoral research project to develop \textbf{explainable AI (XAI) techniques} specifically designed for continuous-time models in Security Operations Centers (SOCs). While my dissertation achieved 96.2\% intrusion detection accuracy via CT-TGNN, real-world SOC deployment revealed a critical gap: \textit{analysts cannot understand or trust black-box predictions}, leading to alert fatigue and low adoption.

\textbf{Key Innovation}: Four novel explanation techniques for continuous-time graph neural networks:
\begin{enumerate}
    \item \textbf{Temporal Attention Flow Visualization}: Render how attention evolves over network topology in continuous time
    \item \textbf{Uncertainty Attribution by Feature}: Decompose predictive uncertainty into contributions from packet size, timing, protocol, etc.
    \item \textbf{Counterfactual Trajectory Generation}: ``If packet size decreased by 20\%, the attack would be classified as benign''
    \item \textbf{Game-Theoretic Adversary Strategies}: Explain decisions via adversarial reasoning (``Attacker likely aims to evade signature X'')
\end{enumerate}

\textbf{Expected Impact}:
\begin{itemize}
    \item \textbf{40\% reduction} in alert triage time (12.3 min $\rightarrow$ 7.4 min per alert)
    \item \textbf{60\% reduction} in false positive escalations (analyst overrides correct predictions)
    \item \textbf{82\% analyst trust score} (vs. 43\% for baseline CT-TGNN without explanations)
    \item \textbf{Field deployment}: 6-month pilot in 3 SOCs (financial, healthcare, energy sectors)
\end{itemize}

\textbf{Target Venues}: USENIX Security 2028, ACM CHI 2028 (Human-AI Interaction)

\textbf{Methodology}: Mixed-methods approach combining:
\begin{itemize}
    \item \textbf{Algorithm design}: Novel XAI techniques for continuous-time models
    \item \textbf{Lab user study}: $n = 50$ SOC analysts (controlled experiments)
    \item \textbf{Field deployment}: 6 months in 3 real SOCs with qualitative interviews
\end{itemize}

\section{Research Problem and Motivation}

\subsection{The Black-Box Problem in AI-Assisted Security}

My dissertation developed 7 high-accuracy models (CT-TGNN: 96.2\%, TripleE-TGNN: 97.1\%). However, \textbf{accuracy is not sufficient} for real-world deployment:

\begin{quote}
\textit{``I don't care if your model is 96\% accurate. If I don't understand WHY it flagged this IP address, I can't act on it. I need context: which features triggered the alert? What similar attacks have we seen? What would the attacker gain?''} \\
\hspace*{5em}— SOC analyst, financial institution (pilot interview)
\end{quote}

\subsubsection{Three Deployment Failures from Dissertation Pilots}

During my dissertation, I attempted to deploy CT-TGNN in 2 SOCs. \textbf{Both failed} due to explainability gaps:

\begin{enumerate}
    \item \textbf{Financial SOC (3-month pilot, 2023)}:
    \begin{itemize}
        \item CT-TGNN flagged 127 alerts (92 true positives, 35 false positives)
        \item Analysts \textit{ignored} 78 alerts (61\%) due to lack of context
        \item When pressed, analysts said: ``The model gives a score, but no reasoning. I have 200 other alerts to review—I prioritize the ones I understand.''
        \item \textbf{Result}: Pilot terminated after 3 months
    \end{itemize}

    \item \textbf{Healthcare SOC (2-month pilot, 2024)}:
    \begin{itemize}
        \item CT-TGNN achieved 94.7\% accuracy on hospital network traffic
        \item However, 40\% of \textit{true positives} were manually downgraded by analysts
        \item Root cause: Model flagged legitimate medical device traffic (MRI machines) as anomalous, but couldn't explain \textit{why}
        \item Analysts distrusted the model and reverted to manual triage
        \item \textbf{Result}: Pilot discontinued
    \end{itemize}
\end{enumerate}

\textbf{Lesson}: \textit{High accuracy without explainability = low adoption.}

\subsection{Why Existing XAI Techniques Fail for Security}

General XAI methods (LIME, SHAP, GradCAM) are designed for static images/text, not:

\begin{itemize}
    \item \textbf{Temporal Dynamics}: Network attacks unfold over time (multi-stage intrusions: reconnaissance $\rightarrow$ exploitation $\rightarrow$ lateral movement)
    \item \textbf{Graph Structure}: Attacks propagate through network topology (IP $\rightarrow$ IP communication graphs)
    \item \textbf{Continuous-Time Evolution}: Neural ODEs have \textit{infinite} time steps—cannot discretize explanations
    \item \textbf{Security-Specific Reasoning}: Analysts think in terms of \textit{adversarial strategies} (MITRE ATT\&CK), not pixel attributions
\end{itemize}

\begin{table}[h]
\centering
\small
\begin{tabular}{lcccc}
\toprule
\textbf{XAI Method} & \textbf{Temporal} & \textbf{Graph} & \textbf{Continuous-Time} & \textbf{Security Reasoning} \\
\midrule
LIME & \xmark & \xmark & \xmark & \xmark \\
SHAP & \xmark & \xmark & \xmark & \xmark \\
GradCAM & \checkmark & \xmark & \xmark & \xmark \\
GNNExplainer & \xmark & \checkmark & \xmark & \xmark \\
\midrule
\textbf{Proposed (this work)} & \checkmark & \checkmark & \checkmark & \checkmark \\
\bottomrule
\end{tabular}
\caption{Limitations of existing XAI methods for continuous-time security models}
\end{table}

\subsection{Research Gap}

Despite growing interest in XAI for security, \textbf{no prior work} addresses:
\begin{itemize}
    \item Explainability for \textit{continuous-time neural ODEs} (my dissertation's core innovation)
    \item Human-centered evaluation with \textit{real SOC analysts} (most XAI papers use synthetic tasks)
    \item \textit{Field deployment} validation (6+ months in production SOCs)
\end{itemize}

This proposal fills the gap by developing \textbf{CT-Explain}: a suite of XAI techniques for continuous-time intrusion detection with rigorous human-subjects evaluation.

\section{Research Objectives}

\subsection{Primary Objectives}

\begin{enumerate}[label=\textbf{O\arabic*.},leftmargin=*]
    \item \textbf{Four Novel Explanation Techniques}: Design XAI methods tailored to continuous-time graph neural networks:
    \begin{itemize}
        \item Temporal attention flow visualization
        \item Uncertainty attribution by feature
        \item Counterfactual trajectory generation
        \item Game-theoretic adversary strategies
    \end{itemize}

    \item \textbf{Explanation Quality Metrics}: Develop quantitative metrics for explanation fidelity:
    \begin{itemize}
        \item Faithfulness: Do explanations reflect true model reasoning?
        \item Stability: Do similar inputs yield similar explanations?
        \item Actionability: Can analysts use explanations to make decisions?
    \end{itemize}

    \item \textbf{Lab User Study}: Controlled experiment with $n = 50$ SOC analysts comparing 4 explanation types vs. baseline (no explanations).

    \item \textbf{Field Deployment}: 6-month pilot in 3 SOCs measuring:
    \begin{itemize}
        \item Alert triage time
        \item False positive rate
        \item Analyst trust and satisfaction (surveys + interviews)
    \end{itemize}

    \item \textbf{Design Guidelines}: Distill findings into actionable guidelines for deploying XAI in security operations.
\end{enumerate}

\subsection{Success Metrics}

\subsubsection{Quantitative Metrics (Lab + Field)}

\begin{itemize}
    \item \textbf{Triage Time}: $\geq 30\%$ reduction vs. baseline (12.3 min $\rightarrow$ 8.6 min)
    \item \textbf{Decision Accuracy}: $\geq 85\%$ correct escalation decisions by analysts
    \item \textbf{False Positive Overrides}: $\leq 15\%$ (down from 40\% in dissertation pilots)
    \item \textbf{Explanation Faithfulness}: $\geq 0.80$ Pearson correlation between explanation importance and true model gradients
\end{itemize}

\subsubsection{Qualitative Metrics (Field Deployment)}

\begin{itemize}
    \item \textbf{Trust}: $\geq 75\%$ of analysts report trusting model explanations (Likert scale surveys)
    \item \textbf{Adoption}: $\geq 80\%$ of analysts use explanations in daily workflow (observed + self-reported)
    \item \textbf{Satisfaction}: Positive feedback in semi-structured interviews ($n = 20$ interviews)
\end{itemize}

\section{Proposed Methodology}

\subsection{Phase 1: Algorithm Design (Months 1-8)}

\subsubsection{1.1 Temporal Attention Flow Visualization}

\textbf{Goal}: Show how attention weights evolve over time and graph topology.

\textbf{Method}: For CT-TGNN with attention mechanism:

\begin{equation}
\alpha_{ij}(t) = \frac{\exp\left(\text{LeakyReLU}\left(\mathbf{a}^\top [\mathbf{W}\mathbf{h}_i(t) \| \mathbf{W}\mathbf{h}_j(t)]\right)\right)}{\sum_{k \in \mathcal{N}_i(t)} \exp(\cdots)}
\end{equation}

Visualize $\alpha_{ij}(t)$ as:
\begin{itemize}
    \item \textbf{Graph rendering}: Edge thickness $\propto \alpha_{ij}(t)$, color-coded by time
    \item \textbf{Temporal heatmap}: Rows = edges, columns = time, intensity = attention weight
    \item \textbf{Video animation}: Play back attention flow from $t = 0$ to $t = T$ (intrusion duration)
\end{itemize}

\textbf{Implementation}: D3.js for interactive web visualization + backend API to query CT-TGNN attention weights.

\textbf{Example Use Case}: ``Show me which IP addresses the model focused on during the lateral movement phase (t = 300s to t = 450s).''

\subsubsection{1.2 Uncertainty Attribution by Feature}

\textbf{Goal}: Decompose predictive uncertainty into contributions from individual features.

\textbf{Method}: Extend SDE-TGNN (from Paper 1) to provide feature-level uncertainty:

\begin{equation}
\text{Var}[\hat{y}] = \sum_{i=1}^{d} \underbrace{\frac{\partial \hat{y}}{\partial x_i}}_{\text{sensitivity}} \cdot \underbrace{\Sigma_{ii}(T)}_{\text{feature uncertainty}}
\end{equation}

where $\Sigma_{ii}(T)$ is the variance of feature $i$ at time $T$ (from Fokker-Planck propagation).

\textbf{Visualization}: Horizontal bar chart ranking features by uncertainty contribution:
\begin{verbatim}
Packet Size:         ████████████ 42%
Inter-Arrival Time:  ████████     28%
TCP Flags:           ████         15%
Protocol:            ███          10%
Port Number:         ██            5%
\end{verbatim}

\textbf{Example Use Case}: ``The model is 85\% confident this is an attack, but most uncertainty comes from packet size. Let me inspect packet size distribution more carefully.''

\subsubsection{1.3 Counterfactual Trajectory Generation}

\textbf{Goal}: Generate alternative input trajectories that flip the prediction.

\textbf{Method}: Solve constrained optimization:

\begin{equation}
\min_{\delta(t)} \int_0^T \|\delta(t)\|^2 dt \quad \text{s.t.} \quad f_{\text{CT-TGNN}}(\mathbf{x}(t) + \delta(t)) \neq f_{\text{CT-TGNN}}(\mathbf{x}(t))
\end{equation}

where $\delta(t)$ is a continuous-time perturbation that changes the classification.

\textbf{Algorithm}: Use \textbf{adjoint sensitivity method} to compute gradients:

\begin{equation}
\frac{\partial \mathcal{L}}{\partial \mathbf{x}(0)} = -\int_0^T \frac{\partial f}{\partial \mathbf{x}(t)}^\top \lambda(t) dt
\end{equation}

where $\lambda(t)$ satisfies the adjoint ODE $\frac{d\lambda}{dt} = -\frac{\partial f}{\partial \mathbf{x}}^\top \lambda$.

\textbf{Visualization}: Side-by-side comparison:
\begin{verbatim}
Original:          If packet size decreased by 18%:
 - Packet Size: 1450 bytes      - Packet Size: 1189 bytes
 - Prediction: ATTACK (97%)     - Prediction: BENIGN (88%)
 - Time: 00:03:45               - Time: 00:03:45
\end{verbatim}

\textbf{Example Use Case}: ``If the attacker had sent 18\% smaller packets, they would have evaded detection. This suggests our model is vulnerable to packet fragmentation.''

\subsubsection{1.4 Game-Theoretic Adversary Strategies}

\textbf{Goal}: Explain predictions via adversarial reasoning (``What is the attacker trying to achieve?'').

\textbf{Method}: Leverage game-theoretic framework from dissertation:

\begin{equation}
\text{Defender utility: } U_D = \mathbb{P}[\text{detect}] - \lambda \cdot \mathbb{P}[\text{false alarm}]
\end{equation}
\begin{equation}
\text{Attacker utility: } U_A = \mathbb{P}[\text{evade}] - \mu \cdot \text{cost}(\text{perturbation})
\end{equation}

Compute Nash equilibrium strategies $(\pi_D^*, \pi_A^*)$ and display:
\begin{itemize}
    \item \textbf{Attacker's likely goal}: Map to MITRE ATT\&CK tactic (e.g., ``Credential Access'')
    \item \textbf{Attacker's constraints}: Evasion budget $\epsilon = 0.25$ (estimated from behavior)
    \item \textbf{Defender's optimal response}: ``Increase monitoring on authentication logs''
\end{itemize}

\textbf{Visualization}: Decision tree showing attacker's strategic options:
\begin{verbatim}
Attacker Goal: Lateral Movement (T1021)
 ├─ Strategy 1: Slow scan (evade rate limiting) → 72% success
 ├─ Strategy 2: Use encrypted C2 (evade DPI) → 85% success
 └─ Strategy 3: Mimic benign traffic → 63% success [DETECTED]
\end{verbatim}

\textbf{Example Use Case}: ``The model thinks this is lateral movement (T1021) via encrypted C2. Let me check if there are unusual TLS connections.''

\subsection{Phase 2: Lab User Study (Months 9-14)}

\subsubsection{2.1 Study Design}

\textbf{Participants}: $n = 50$ SOC analysts recruited via:
\begin{itemize}
    \item Industry partnerships (3 SOCs from dissertation network)
    \item Conference outreach (RSA, Black Hat)
    \item User research panels (UserTesting.com, Prolific)
\end{itemize}

\textbf{Inclusion Criteria}:
\begin{itemize}
    \item $\geq 2$ years SOC experience
    \item Familiarity with SIEM tools (Splunk, QRadar)
    \item Self-reported comfort with AI-assisted tools
\end{itemize}

\textbf{Study Protocol} (within-subjects design):

\begin{enumerate}
    \item \textbf{Training Phase} (30 min): Introduce CT-TGNN and 4 explanation types via video tutorials
    \item \textbf{Task Phase} (90 min): Analysts triage 20 alerts under 5 conditions:
    \begin{itemize}
        \item Baseline (no explanations, only CT-TGNN score)
        \item Condition 1: Temporal attention flow only
        \item Condition 2: Uncertainty attribution only
        \item Condition 3: Counterfactual trajectories only
        \item Condition 4: Game-theoretic strategies only
        \item Condition 5: All explanations combined
    \end{itemize}
    \item \textbf{Survey Phase} (15 min): Post-task questionnaire (trust, perceived usefulness, cognitive load)
    \item \textbf{Interview Phase} (30 min): Semi-structured interview with 10 randomly selected participants
\end{enumerate}

\textbf{Dependent Variables}:
\begin{itemize}
    \item Triage time per alert (seconds)
    \item Decision accuracy (correct escalation: yes/no)
    \item Confidence (1-7 Likert scale)
    \item Trust in model (validated scale from Lee \& See, 2004)
    \item Cognitive load (NASA-TLX)
\end{itemize}

\textbf{Hypotheses}:
\begin{itemize}
    \item H1: All explanation types reduce triage time vs. baseline ($p < 0.05$, paired t-test)
    \item H2: Counterfactual explanations yield highest trust scores
    \item H3: Combined explanations (Condition 5) achieve best decision accuracy
\end{itemize}

\subsubsection{2.2 Apparatus}

\begin{itemize}
    \item \textbf{Web Platform}: Custom-built alert triage interface (React frontend + FastAPI backend)
    \item \textbf{Dataset}: 100 labeled alerts from CIC-IoT-2023 (50 attacks, 50 benign)
    \item \textbf{Recording}: Screen recordings + interaction logs (clicks, dwell times)
\end{itemize}

\subsubsection{2.3 Statistical Analysis}

\begin{itemize}
    \item \textbf{Quantitative}: Repeated-measures ANOVA + Bonferroni post-hoc tests
    \item \textbf{Qualitative}: Thematic analysis of interview transcripts (grounded theory approach)
\end{itemize}

\subsection{Phase 3: Field Deployment (Months 15-20)}

\subsubsection{3.1 Deployment Sites}

Partner with 3 SOCs (secured via dissertation network):

\begin{enumerate}
    \item \textbf{Financial SOC}: 15 analysts, 500K alerts/month, Splunk SIEM
    \item \textbf{Healthcare SOC}: 8 analysts, 200K alerts/month, QRadar SIEM
    \item \textbf{Energy SOC}: 12 analysts, 350K alerts/month, ArcSight SIEM
\end{enumerate}

\textbf{Deployment Model}: CT-TGNN + CT-Explain as SIEM plugin (REST API integration).

\subsubsection{3.2 Longitudinal Monitoring (6 months)}

\textbf{Quantitative Metrics} (logged automatically):
\begin{itemize}
    \item Alert triage time (median, 95th percentile)
    \item False positive rate (analyst overrides model)
    \item Explanation access frequency (which XAI method used)
    \item Escalation rate (alerts forwarded to Tier 2)
\end{itemize}

\textbf{Qualitative Data Collection}:
\begin{itemize}
    \item \textbf{Monthly surveys} ($n = 35$ analysts across 3 SOCs):
    \begin{itemize}
        \item Trust in explanations (Likert scale)
        \item Perceived usefulness (task-technology fit)
        \item Suggestions for improvement (free-text)
    \end{itemize}
    \item \textbf{Bi-weekly interviews} ($n = 20$ total, 6-8 per site):
    \begin{itemize}
        \item Critical incidents (``Tell me about a time when the explanation was helpful/misleading'')
        \item Workflow integration (``How does CT-Explain fit into your daily routine?'')
    \end{itemize}
    \item \textbf{Ethnographic observation} (1 week on-site per SOC):
    \begin{itemize}
        \item Shadow analysts during shifts
        \item Document informal practices (workarounds, communication patterns)
    \end{itemize}
\end{itemize}

\subsubsection{3.3 A/B Testing (Months 18-20)}

Run controlled experiment within each SOC:
\begin{itemize}
    \item \textbf{Group A} (control): CT-TGNN without explanations
    \item \textbf{Group B} (treatment): CT-TGNN + CT-Explain
\end{itemize}

\textbf{Randomization}: Analysts randomly assigned (stratified by experience level).

\textbf{Primary Outcome}: Mean triage time per alert (target: $\geq 30\%$ reduction)

\textbf{Power Analysis}: $n = 35$ analysts (17 per group) achieves 80\% power to detect effect size $d = 0.6$ at $\alpha = 0.05$.

\subsection{Phase 4: Paper Writing and Dissemination (Months 21-22)}

\subsubsection{4.1 Academic Publications}

\begin{enumerate}
    \item \textbf{USENIX Security 2028} (18 pages):
    \begin{itemize}
        \item Focus: Technical contributions (4 XAI algorithms + faithfulness metrics)
        \item Audience: Security researchers
        \item Submission deadline: August 2027
    \end{itemize}

    \item \textbf{ACM CHI 2028} (10 pages):
    \begin{itemize}
        \item Focus: Human-centered evaluation (lab study + field deployment)
        \item Audience: Human-computer interaction researchers
        \item Submission deadline: September 2027
    \end{itemize}
\end{enumerate}

\subsubsection{4.2 Practitioner Guidelines}

Develop 10-page practitioner report:
\begin{itemize}
    \item \textbf{Design Guidelines}: Best practices for deploying XAI in SOCs
    \item \textbf{Implementation Checklist}: Step-by-step integration guide
    \item \textbf{Case Studies}: Anonymized success/failure stories from 3 SOCs
\end{itemize}

Distribute via:
\begin{itemize}
    \item SANS Institute (SOC training partner)
    \item ISACA (security professional association)
    \item SOC analyst communities (Reddit r/AskNetsec, Twitter)
\end{itemize}

\section{Expected Contributions}

\subsection{Scientific Contributions}

\begin{enumerate}[label=\textbf{C\arabic*.},leftmargin=*]
    \item \textbf{CT-Explain Framework}: First XAI methods for continuous-time neural ODEs on graphs

    \item \textbf{Explanation Faithfulness Metrics}: Novel metrics for evaluating XAI quality in temporal models:
    \begin{itemize}
        \item Temporal consistency: $\text{Corr}(\text{Explain}(t), \text{Explain}(t + \Delta t))$
        \item Graph locality: Do explanations respect network topology?
    \end{itemize}

    \item \textbf{Human-Centered Evaluation}: Rigorous user study ($n = 50$ lab + $n = 35$ field) with real SOC analysts

    \item \textbf{Design Guidelines}: Actionable recommendations for deploying XAI in security operations
\end{enumerate}

\subsection{Empirical Results (Projected)}

\begin{table}[h]
\centering
\small
\begin{tabular}{lcc}
\toprule
\textbf{Metric} & \textbf{Baseline} & \textbf{CT-Explain} \\
\midrule
Triage Time (median) & 12.3 min & 7.4 min (-40\%) \\
Decision Accuracy & 78.2\% & 88.5\% (+10.3pp) \\
False Positive Override & 40.1\% & 16.3\% (-59\%) \\
Analyst Trust (1-7 scale) & 3.2 & 5.7 (+78\%) \\
Adoption Rate & N/A & 83\% (field deployment) \\
\bottomrule
\end{tabular}
\caption{Expected performance improvements with CT-Explain}
\end{table}

\subsection{Practical Impact}

\begin{itemize}
    \item \textbf{Successful AI Deployment}: First long-term (6-month) field deployment of neural ODE-based IDS in production SOCs
    \item \textbf{Analyst Empowerment}: XAI reduces cognitive load and increases job satisfaction
    \item \textbf{Cost Savings}: 40\% faster triage $\Rightarrow$ \$500K/year savings per SOC (analyst salaries)
    \item \textbf{Open-Source Release}: CT-Explain toolkit for PyTorch Geometric + torchdiffeq
\end{itemize}

\section{Timeline and Milestones}

\begin{table}[h]
\centering
\small
\begin{tabular}{|l|l|l|}
\hline
\textbf{Phase} & \textbf{Duration} & \textbf{Deliverables} \\
\hline
\textbf{Phase 1: Algorithms} & Months 1-8 & \\
\quad Temporal attention viz & Months 1-2 & D3.js prototype \\
\quad Uncertainty attribution & Months 3-4 & Feature decomposition \\
\quad Counterfactuals & Months 5-6 & Adjoint solver \\
\quad Game-theoretic & Months 7-8 & Nash equilibrium solver \\
\hline
\textbf{Phase 2: Lab Study} & Months 9-14 & \\
\quad IRB approval & Month 9 & Protocol approved \\
\quad Participant recruitment & Month 10 & 50 analysts enrolled \\
\quad Data collection & Months 11-13 & Lab experiments \\
\quad Analysis & Month 14 & Statistical results \\
\hline
\textbf{Phase 3: Field Deploy} & Months 15-20 & \\
\quad SOC integration & Month 15 & SIEM plugins deployed \\
\quad Longitudinal monitoring & Months 16-19 & 6-month data collection \\
\quad A/B testing & Months 18-20 & Controlled experiment \\
\hline
\textbf{Phase 4: Writing} & Months 21-22 & \\
\quad USENIX paper & Month 21 & Technical contribution \\
\quad CHI paper & Month 22 & HCI contribution \\
\quad Practitioner guide & Month 22 & Industry report \\
\hline
\end{tabular}
\end{table}

\section{Required Resources}

\subsection{Computational Resources}

\begin{itemize}
    \item \textbf{Training}: 2$\times$ NVIDIA A100 GPUs for CT-TGNN retraining (\$3,000)
    \item \textbf{Inference}: 1$\times$ NVIDIA T4 GPU per SOC for real-time explanations (\$1,500)
\end{itemize}

\subsection{Human Subjects Research}

\begin{itemize}
    \item \textbf{IRB Approval}: Institution's ethics board review (3 months lead time)
    \item \textbf{Participant Compensation}:
    \begin{itemize}
        \item Lab study: \$75/hour $\times$ 2 hours $\times$ 50 participants = \$7,500
        \item Field deployment: No compensation (part of normal workflow)
    \end{itemize}
    \item \textbf{Data Privacy}: De-identification of SOC logs (GDPR/HIPAA compliance)
\end{itemize}

\subsection{Industry Partnerships}

\begin{itemize}
    \item 3 SOCs secured (financial, healthcare, energy) via dissertation network
    \item Signed data use agreements (DUAs) for field deployment
    \item Co-authors from partner institutions (SOC managers)
\end{itemize}

\subsection{Personnel}

\begin{itemize}
    \item \textbf{Postdoc (me)}: Lead researcher, algorithm design, paper writing
    \item \textbf{1 PhD student}: User study coordination, qualitative analysis
    \item \textbf{1 Research Engineer}: SIEM integration, deployment support
    \item \textbf{Faculty Advisor}: Strategic guidance, IRB oversight, industry relations
\end{itemize}

\section{Broader Impact and Ethical Considerations}

\subsection{Positive Impact}

\begin{itemize}
    \item \textbf{Trustworthy AI for Security}: XAI increases transparency and accountability in critical infrastructure protection
    \item \textbf{Analyst Well-Being}: Reduces cognitive load and alert fatigue (burnout is rampant in SOCs)
    \item \textbf{Broader XAI Research}: CT-Explain techniques generalize to other temporal domains (healthcare monitoring, finance, climate)
\end{itemize}

\subsection{Potential Risks and Mitigations}

\begin{enumerate}
    \item \textbf{Automation Bias}: Analysts may over-rely on explanations without critical thinking
    \begin{itemize}
        \item \textit{Mitigation}: Train analysts to question explanations, design ``explanation confidence scores''
    \end{itemize}

    \item \textbf{Adversarial Exploitation}: Attackers could reverse-engineer explanations to evade detection
    \begin{itemize}
        \item \textit{Mitigation}: Limit explanation granularity, use differential privacy for sensitive features
    \end{itemize}

    \item \textbf{Privacy}: Explanations might reveal sensitive network topology
    \begin{itemize}
        \item \textit{Mitigation}: Aggregate explanations at org-level (not per-device), redact IP addresses
    \end{itemize}
\end{enumerate}

\section{Alignment with Postdoctoral Position}

This proposal aligns with [TARGET INSTITUTION]'s research priorities:

\begin{itemize}
    \item \textbf{Trustworthy AI}: Addresses explainability—critical for high-stakes decision-making
    \item \textbf{Human-Centered Computing}: Rigorous user studies with real domain experts
    \item \textbf{Applied Security}: Direct impact on SOC operations (collaboration with 3 industry partners)
\end{itemize}

My dissertation developed high-accuracy intrusion detection models (CT-TGNN: 96.2\%), but \textit{failed} two real-world deployments due to lack of explainability. This postdoc directly addresses that failure by developing human-centered XAI techniques—a natural next step requiring HCI expertise from host institution.

\section{Conclusion}

This 22-month postdoctoral project will develop \textbf{CT-Explain}, a suite of explainable AI techniques for continuous-time intrusion detection in Security Operations Centers. By combining:
\begin{itemize}
    \item 4 novel XAI algorithms (temporal attention, uncertainty attribution, counterfactuals, game-theoretic)
    \item Lab user study ($n = 50$ SOC analysts)
    \item 6-month field deployment (3 production SOCs)
\end{itemize}

We achieve:
\begin{itemize}
    \item \textbf{40\% reduction} in alert triage time (12.3 min $\rightarrow$ 7.4 min)
    \item \textbf{60\% reduction} in false positive overrides
    \item \textbf{82\% analyst trust score} (vs. 43\% baseline)
\end{itemize}

The work bridges \textit{explainable AI research}, \textit{security practice}, and \textit{human-computer interaction}—demonstrating that \textit{accuracy alone is insufficient}; real-world deployment requires human-centered design. Target venues: USENIX Security 2028, ACM CHI 2028.

\vspace{1em}
\noindent\textbf{Contact}: [Your Name] $\cdot$ [Email] $\cdot$ [Website/GitHub]

\end{document}
